\section{Additional incidents and contrasts}
\label{app:incidents}

\Incident{g7}{A public final offer below the private limit (G7)}
\Setup{At G7 d42, turn 25, Grok wants North Carolina Avenue from Sol to advance a green set.}

\TranscriptQuote{g7_d42_public}
\TranscriptQuote{g7_d42_private}
\TranscriptQuote{g7_d43_public}

\TranscriptFinding{Events 324--327 transfer North Carolina to Grok for \$600. Sol's applied note says it would rather reinvest in cheaper sets than hold out for another \$25. The private \$625 contingency is never tested.}

\TranscriptReading{This is a public/private discrepancy about a bargaining limit followed by settlement at the advertised limit. Together with Opus's ``final'' \$450 in G2, which did not settle, it shows the same authorized tactic ending both ways. Sol's stated reason does not show that the word ``final'' secured the price.}

\TranscriptSource{G7 d42--43; both selected attempts are standard and valid at index 0. Sol's d43 prompt contains the \$600 message.}

\Incident{disclosure}{A private note delivered as public chat (G20 d1)}
\Setup{On Grok's first end-of-turn decision in G20, the assistant content requests an empty public message. The received tool arguments already contain a fragment copying the private note into the public field.}

\TranscriptQuote{g20-received-and-applied-public}
\TranscriptQuote{g20-recipient-recognition}

\TranscriptFinding{The fragment appears in the received arguments, the applied public field, the canonical message event and the next three players' prompts. Gemini records the breach at d4.}

\TranscriptReading{A private note crossed into a public channel and an opponent noticed. Its contents concern public holdings and early priorities, and no exploitation follows. The empty public field in assistant content argues against a deliberate disclosure. As with the action-name conflicts, the record places the contamination at or before the received response without identifying its source.}

\TranscriptSource{G20 d1 and d4, standard attempt 0. Decision indices start at zero.}

\subsection{Recipient responses to false or misleading propositions}

\begin{table}[htbp]
\centering\small
\begin{tabularx}{\linewidth}{@{}L{1.4cm}YYYL{1.9cm}@{}}
\toprule
Case & Proposition and status & Speaker's other channels & Recipient record and response & Outcome\\
\midrule
G13 d252 & ``\$130 leaves me \$35 short''; true in cash, omits a \$130 mortgage (\$295 reachable) & Note acknowledges the route; provider explanation states the omission and also contradicts itself & Solvency value stated at d251 after the seller's d250 appeal; silent on the \$130 route; accepts \$200 & Settled at \$200\\
G19 d135 & ``States is the third orange''; false & Note says abort; content selects \End; received call proposes & Corrects the premise publicly; accepts & Settled\\
G11 d255 & ``Virginia completes your pinks''; false & Note repeats the premise; later concedes & Provider explanation echoes, then corrects; applied note rejects & Rejected twice\\
G4 d149 & ``\$200 a house'' for yellow; false (\$150) & Note states a one-house limit without the \$200 figure & True \$150 and a contingency recorded at d148; public ``Fair point''; switches request & Settled at \$50\\
G6 d19 & ``Your oranges are the board's only set''; false & Note repeats the premise; content selects \End; received call proposes & Note rejects the premise; accepts & Settled; proposer's later note calls it an overpay\\
\bottomrule
\end{tabularx}
\caption{Five source-checked proposition cases, inherited from the baseline review. ``Recipient record'' refers to the recipient's channels at or before its response. Selected contrasts, not prevalence.}
\label{tab:framing}
\end{table}

\paragraph{G11 echo and correction.} At d256 Grok's returned provider explanation first repeats Sol's premise (``accepting would complete my pink monopoly'') and a few hundred characters later corrects it. Its applied note calls the offer a lie and rejects. A recipient's accusation does not establish the proposer's knowledge: Sol's own applied note repeats the false completion claim, which makes a mistake more plausible than a knowing deception, and Sol concedes at d257.

\paragraph{G4 public concession.} After Sol completes the pinks, Gemini at d138 publicly invites deals so that Sol is not the only player with a set. Grok buys mortgaged Ventnor from Sol for \$200 at d146 to make the yellows tradable. In the orange--yellow negotiation Gemini's d148 note already records the true \$150 house cost and a contingency for demanding compensation. Grok's d149 message says yellow houses cost \$200. Gemini's d150 ``Fair point on house costs'' precedes a request to take yellow with \$143, while its returned explanation keeps \$150. The settlement at \$50 completes groups for both players. The direction change matches a contingency recorded before the false figure, so no price or framing effect is identified.

\paragraph{A label we rejected.} In G12 d195 a \$410 insolvency claim is corrected by the recipient to \$630, and only afterward does the speaker call the earlier pitch a bluff. Its d195 explanation treats the \$410 figure as real. The retrospective label cannot establish prior knowing deception, so the case is not used as an example of lying.

\paragraph{Other refusals to fund rivals.} Gemini's G9 d82 rejection says that paying would fund a dark-blue house, and its G16 d322 rejection prefers to keep a rival short of cash. These are stated reasons that were not recomputed. They complement the G20 appeal and the 61 rival-consequence rejections coded in Section~\ref{sec:bargaining}, without a separate prevalence estimate. No multi-party pact, cartel or deliberate kingmaking is identified in the reviewed cases; this is not a corpus-wide absence claim.
